\documentclass[9pt]{extarticle}

\usepackage[utf8]{inputenc}
\usepackage[a4paper, margin=1.2cm]{geometry}
\usepackage[colorlinks=true, urlcolor=blue]{hyperref}
\usepackage{enumitem}
\usepackage{changepage}

\pagestyle{empty}

\setlist[itemize]{label=\textbullet, noitemsep, topsep=0pt}
\setlength{\parindent}{0pt}

\begin{document}

\begin{center}
	{\huge\bfseries Paul Parisot}
	\\[4pt]
	{\Large Ingénieur Logiciel - C++, Lua \& Systèmes Temps Réel}
	\\
\end{center}

\par\vspace{-5pt}
\rule{\linewidth}{0.4pt}
\par\vspace{-5pt}

\begin{center}
	Github: \href{https://github.com/paulogarithm}{paulogarithm}
	\\
	\href{mailto:paul.parisot@epitech.eu}{paul.parisot@epitech.eu}
	\\
\end{center}

\par\vspace{-10pt}
\rule{\linewidth}{0.4pt}
\par\vspace{7pt}

{\large\bfseries À propos}
\par\vspace{-8pt}
\rule{\linewidth}{0.4pt}
\par\vspace{5pt}

Ingénieur logiciel spécialisé en C/C++ et scripting embarqué (Lua, Python), avec un intérêt pour les {\bfseries systèmes temps réel, le contrôle de systèmes et l'automatisation}. Attentif à la fiabilité du code : tests, Valgrind, CI/CD.
\\
Expérience concrète d'intégration Lua dans des moteurs de jeu C++ et des serveurs réseau, et développement de modules Luau (génération procédurale, algorithmes) sur Roblox Studio.
\\
Un an d'études à Montréal (McGill University, management) et pratique professionnelle de l'anglais (C1). En préparation de l'agrégation d'informatique.
\par\vspace{5pt}

{\large\bfseries Compétences}
\par\vspace{-8pt}
\rule{\linewidth}{0.4pt}
\par\vspace{5pt}

\hspace*{5pt}{\bfseries Maîtrise}: C/C++, Lua/Luau, Python, Golang, Assembleur, Shell, SQL
\\
\hspace*{5pt}{\bfseries Scripting \& Embarqué}: intégration Lua dans C/C++, modules Luau (algorithmes, génération procédurale), Arduino (notions)
\\
\hspace*{5pt}{\bfseries Outils}: gdb, Valgrind, Ghidra, Git, Docker, CI/CD, Environnement Linux
\\
\hspace*{5pt}{\bfseries Langues}: Français, Anglais (C1, compétences professionnelles)

\par\vspace{5pt}

{\large\bfseries Expérience Professionnelle}
\par\vspace{-8pt}
\rule{\linewidth}{0.4pt}
\par\vspace{5pt}

{\bfseries Gno.land - Ingénieur Logiciel Junior (Open-Source)}
\begin{itemize}[noitemsep, topsep=0pt]
	\item Développement de 6+ packages core en Gnolang (variante de Go pour smart contracts), dont une compatibilité EVM (traduction d'opcodes en instructions Go).
	\item Détection et correction de {\bfseries 6 bugs critiques} causant des crashs runtime ; {\bfseries ~10k lignes de code} contribuées.
	\item Rédaction de tests exhaustifs sous des politiques strictes de CI/CD et de revue de code.
\end{itemize}

\par\vspace{5pt}

{\bfseries Faurecia / Forvia - Ingénieur Démo Logicielle}
\begin{itemize}[noitemsep, topsep=0pt]
	\item Correction d'un ancien code {\bfseries C++} pour le faire communiquer avec les bons drivers matériels (LED) ; scripts Python de pilotage de hardware.
	\item Conception d'un {\bfseries moteur de calcul d'itinéraire GPS} et d'une interface de recherche de POI, intégrés à une application démo présentée au {\bfseries CES 2024}.
	\item Livraison de fonctionnalités prêtes pour la production {\bfseries en avance sur le planning}.
\end{itemize}

\par\vspace{5pt}

{\bfseries Epitech / IONIS - Assistant Technique \& Formateur}
\begin{itemize}[noitemsep, topsep=0pt]
	\item Enseignement à {\bfseries plus de 400 étudiants} en C, C++, asm et Haskell ; correction et évaluation des rendus.
	\item Création de {\bfseries contenus pédagogiques interactifs} (quiz, cours inversés, tutoriels).
\end{itemize}

\par\vspace{5pt}

{\large\bfseries Projets Personnels}
\par\vspace{-8pt}
\rule{\linewidth}{0.4pt}
\par\vspace{5pt}

\underline{\bfseries \href{https://github.com/paulogarithm/blorox}{BloRox} -- Moteur 3D temps réel en C++ avec scripting Lua}
\par
\begin{adjustwidth}{5pt}{0pt}
	Architecture ECS orientée données : {\bfseries 100k entités animées}, $\sim$6 ms de moteur et $\sim$6 ms de rendu par image. \\
	{\bfseries Intégration d'un moteur de script Lua} pour le comportement des entités et la logique de jeu, avec binding C++/Lua.
\end{adjustwidth}
\par
{\bfseries Compétences :} C++23/26, ECS, Lua, binding C++/Lua, alignement mémoire, SIMD, optimisation de performance

\par\vspace{8pt}

\underline{\bfseries Zappy -- Serveur TCP multi-clients en C avec intégration Lua}
\par
\begin{adjustwidth}{5pt}{0pt}
    Serveur TCP mono-thread à boucle événementielle ({\bfseries \texttt{select}}), avec {\bfseries intégration d'un interpréteur Lua en C} pour la personnalisation du comportement des entités. \\
    Buffer borné par client et mise en quarantaine des clients qui le saturent ; aucune erreur ni fuite sous Valgrind.
\end{adjustwidth}
\par
{\bfseries Compétences :} C, sockets TCP, select, Lua (intégration C/Lua), gestion mémoire, Valgrind

\par\vspace{8pt}

\underline{\bfseries Modules Luau (Roblox Studio) -- Algorithmes et génération procédurale}
\par
\begin{adjustwidth}{5pt}{0pt}
    Développement de {\bfseries modules Luau} pour Roblox Studio : {\bfseries génération aléatoire d'arbres et de maisons} (algorithmes procéduraux), {\bfseries décodeur PNG} en Luau. \\
    Conception de systèmes de jeu avec typage progressif et optimisation des performances.
\end{adjustwidth}
\par
{\bfseries Compétences :} Luau, algorithmes procéduraux, typage, optimisation, Roblox Studio

\par\vspace{8pt}

\underline{\bfseries \href{https://github.com/paulogarithm/cclass}{cclass} -- Moteur de mutation mémoire à l'exécution}
\par
\begin{adjustwidth}{5pt}{0pt}
    Framework en C pur modifiant dynamiquement des instructions assembleur compilées directement en mémoire, pour contourner le surcoût du dispatch par vtable. \\
    Appels système {\bfseries \texttt{mmap}} avec des flags de protection précis et règles d'alignement mémoire strictes.
\end{adjustwidth}
\par
{\bfseries Compétences :} C, systèmes Linux, assembleur x86, mmap, sûreté mémoire

\par\vspace{8pt}

{\large\bfseries Formation}
\par\vspace{-8pt}
\rule{\linewidth}{0.4pt}
\par\vspace{5pt}

{\bfseries Epitech, Paris} \hfill 2022 -- 2027 \\
Diplôme d'ingénieur en technologies de l'information (diplôme de niveau Master) \\
Certification Niveau 7 - Expert en ingénierie logicielle

\vspace{6pt}

{\bfseries McGill University, Montréal} \hfill 2025 -- 2026 \\
Certificat en Management -- GPA : 3.57/4.0 \\
Cours suivis : Comptabilité, Managerial Economics \& Analysis, Business, Mathématiques, Gestion de projet, Marketing

\end{document}
